\chapter{数列の極限}\label{ch:text3}

\begin{textref}{第3章 1節--6節}
  数列の極限の $\varepsilon$--$N$ 論法による定義(定義1.2)と極限の性質(1--2節: 命題1.1, 1.3, 定理2.1--2.18), 等比数列の極限(3節: 命題3.1), 級数(4節: 定義4.1, 命題4.3, 定理4.4, 4.7), 有界な単調数列(5節: 定理5.3, 5.10), 部分列・Bolzano--Weierstrass の定理・Cauchy 列(6節: 定理6.7, 6.8, 補題6.10, 定理6.11)を扱う. 本章の各定理の直後の「形式化」欄に形式的な主張と第III部の証明木の頁を示す.
\end{textref}

\section{極限の定義}

以下, $n, N, N_1, N_2$ は $\omega$ の元を表す変数, $N \leq n$ は第\ref{ch:zfc}章の順序である. 数列は $\Seq(a) :\leftrightarrow a : \N \to \R$ であり, $a_n := a(n)$. 数列について用いる集合論的な事実(値が実数であること, 項ごとの和・定数倍の数列の存在と値)は付録の基礎定理 [S1]--[S5] である.

\begin{dfn}[テキスト定義1.2(1)]\label{dfn:conv}
  述語記号 $\Conv$ を次の定義公理で導入する.
  \begin{align*}
    \Conv(a, \alpha) :\leftrightarrow{}& \Seq(a) \land \alpha \in \R \\
    & \land \forall \varepsilon\,\bigl(\varepsilon \in \R \to 0 < \varepsilon \to \exists N\,(N \in \N \land \forall n\,(n \in \N \to N \leq n \to \abs{a_n - \alpha} < \varepsilon))\bigr).
  \end{align*}
  $\Conv(a, \alpha)$ を「数列 $a$ は $\alpha$ に収束する」と読む. $\mathrm{Cgt}(a) :\leftrightarrow \exists \alpha\,\Conv(a, \alpha)$.
\end{dfn}

\begin{dfn}[テキスト定義1.2(2)]
  \begin{align*}
    \mathrm{Dvg}(a) &:\leftrightarrow \Seq(a) \land \neg\mathrm{Cgt}(a),\\
    \mathrm{DvgP}(a) &:\leftrightarrow \Seq(a) \\
    &\qquad \land \forall K\,\bigl(K \in \R \to 0 < K \to \exists N\,(N \in \N \land \forall n\,(n \in \N \to N \leq n \to K < a_n))\bigr).
  \end{align*}
  テキストの「$\lim a_n = \infty$」は $\mathrm{DvgP}(a)$ の別表記である. $-\infty$ への発散も同様に定める.
\end{dfn}

\begin{lem}\label{lem:conv-unfold}
  \begin{enumerate}[label={[C\arabic*]}]
    \item $\Conv(a, \alpha) \to \Seq(a)$.
    \item $\Conv(a, \alpha) \to \alpha \in \R$.
    \item $\Conv(a, \alpha) \to \varepsilon \in \R \to 0 < \varepsilon \to \exists N\,(N \in \N \land \forall n\,(n \in \N \to N \leq n \to \abs{a_n - \alpha} < \varepsilon))$.
  \end{enumerate}
\end{lem}
\begin{proof}[通常の証明]
  定義公理の例化と命題論理.
\end{proof}

\trees{ConvSeq,ConvR,ConvE}
\section{命題1.1}

\begin{prop}[テキスト命題1.1]\label{prop:eps-eq}
  $\ZFC \vdash \forall a\forall b\,(a \in \R \to b \in \R \to \forall\varepsilon(\varepsilon \in \R \to 0 < \varepsilon \to \abs{a - b} < \varepsilon) \to a = b)$.
\end{prop}
\begin{proof}[通常の証明]
  $a \neq b$ と仮定すると $0 < \abs{a - b}$ ([A15][A14]). 仮定を $\varepsilon := \abs{a - b}$ に適用すると $\abs{a - b} < \abs{a - b}$ となり [O1] に反する.
\end{proof}

\trees{EpsEq}
\section{尾部の共通化}

テキストの定理2.1, 2.3, 2.5 の証明はいずれも「$N(\varepsilon) = \max\{N_1, N_2\}$ とおけば $n \geq N(\varepsilon)$ で両方の評価が成り立つ」という同じ論法を含む. これを1つの補題として証明する. $\max\{N_1, N_2\} = N_1 \cup N_2$ であり, $N_1 \leq N_1 \cup N_2$ は集合一般の事実である([C4]--[C8]).

\begin{lem}\label{lem:cup}
  \begin{enumerate}[label={[C\arabic*]}, start=4]
    \item $m \subseteq m \cup n$.
    \item $n \subseteq m \cup n$.
    \item $x \subseteq y \to x \leq y$.
    \item $m \leq m \cup n$.
    \item $n \leq m \cup n$.
  \end{enumerate}
\end{lem}

\trees{SubCupL,SubCupR,SubLe,LeCupL,LeCupR}
\begin{lem}[尾部の共通化 [C9]]\label{lem:tail}
  \begin{multline*}
    \Conv(a, \alpha) \to \Conv(b, \beta) \to \varepsilon \in \R \to 0 < \varepsilon \\
    \to \exists N\,\bigl(N \in \N \land \forall n\,(n \in \N \to N \leq n \to \abs{a_n - \alpha} < \varepsilon \land \abs{b_n - \beta} < \varepsilon)\bigr).
  \end{multline*}
\end{lem}
\begin{proof}[通常の証明]
  [C3] から $N_1, N_2$ をとり, $M := N_1 \cup N_2$ とおく. $M \in \N$ ([N3]). $n \in \N$, $M \leq n$ なら [C7][C8][O5] より $N_1 \leq n$, $N_2 \leq n$ なので両方の評価が成り立つ.
\end{proof}

\trees{Tail}
\section{極限の一意性}

\begin{lem}[[C10]]
  $x, \alpha, \beta, \varepsilon \in \R$ について $\abs{x - \alpha} < \varepsilon \to \abs{x - \beta} < \varepsilon \to \abs{\alpha - \beta} < \varepsilon + \varepsilon$.
\end{lem}
\begin{proof}[通常の証明]
  $\alpha - \beta = (\alpha - x) + (x - \beta)$ ([F16]) なので, 三角不等式と [A13] より $\abs{\alpha - \beta} \leq \abs{x - \alpha} + \abs{x - \beta} < \varepsilon + \varepsilon$ ([F20], [O4]).
\end{proof}

\trees{ErrTwo}
\begin{prop}[極限の一意性. テキスト命題1.3]\label{thm:lim-unique}
  $\ZFC \vdash \Conv(a, \alpha) \to \Conv(a, \beta) \to \alpha = \beta$.
\end{prop}
\begin{proof}[通常の証明]
  $\alpha \neq \beta$ と仮定し, $\varepsilon := \abs{\alpha - \beta}/2 > 0$ とおく. 補題\ref{lem:tail}($b := a$)より, ある $N$ について $\abs{a_N - \alpha} < \varepsilon$, $\abs{a_N - \beta} < \varepsilon$. [C10] より $\abs{\alpha - \beta} < \varepsilon + \varepsilon = \abs{\alpha - \beta}$ ([F36]) となり [O1] に反する.
\end{proof}

\trees{LimUnique}
\begin{dfn}[$\lim$]
  命題\ref{thm:lim-unique}と定義\ref{dfn:description}により, 関数記号 $\lim a$ を「$\Conv(a, \alpha)$ をみたす唯一の $\alpha$」で定める. テキストの $\lim_{n\to\infty} a_n = \alpha$ は $\Conv(a, \alpha)$ を意味する.
\end{dfn}

\section{和と定数倍の極限}

数列 $a, b$ と $c \in \R$ に対して, 項ごとの和 $a \oplus b := (n \in \N \mapsto a_n + b_n)$ と定数倍 $c \odot a := (n \in \N \mapsto c \cdot a_n)$ を定義する(定義\ref{dfn:lambda}). これらが数列であり値が $a_n + b_n$, $c \cdot a_n$ であることは基礎定理 [S2][S3] である.

\begin{lem}[[C11]]
  $x, y, \alpha, \beta, h \in \R$ について $\abs{x - \alpha} < h \to \abs{y - \beta} < h \to \abs{(x + y) - (\alpha + \beta)} < h + h$.
\end{lem}

\trees{ErrSum}
\begin{thm}[和の極限. テキスト定理2.1(1)]\label{thm:lim-sum}
  $\ZFC \vdash \Conv(a, \alpha) \to \Conv(b, \beta) \to \Conv(a \oplus b, \alpha + \beta)$.
\end{thm}
\begin{proof}[通常の証明(テキストと同じ)]
  任意の $\varepsilon > 0$ をとる. 補題\ref{lem:tail}を $\varepsilon/2 > 0$ ([F35]) に適用して $N$ をとると, $n \geq N$ ならば $\abs{a_n - \alpha} < \varepsilon/2$, $\abs{b_n - \beta} < \varepsilon/2$ であり, [C11] と [F36] より $\abs{(a_n + b_n) - (\alpha + \beta)} < \varepsilon$.
\end{proof}

\trees{LimSum}
\begin{thm}[定数倍の極限. テキスト定理2.1(2)]\label{thm:lim-scalar}
  $\ZFC \vdash c \in \R \to \Conv(a, \alpha) \to \Conv(c \odot a, c\alpha)$.
\end{thm}
\begin{proof}[通常の証明(テキストと同じ)]
  $c = 0$ なら $\abs{0 \cdot a_n - 0 \cdot \alpha} = 0 < \varepsilon$ なので $N = 1_\omega$ でよい. $c \neq 0$ なら $0 < \abs{c}$ ([A14]) であり, $\varepsilon \abs{c}^{-1} > 0$ ([F34], (R15)). これに対する $N$ をとると, $n \geq N$ で $\abs{c a_n - c\alpha} = \abs{c}\abs{a_n - \alpha} < \abs{c} \cdot \varepsilon\abs{c}^{-1} = \varepsilon$ ([F17], [A17], [F32]).
\end{proof}

\trees{LimScalar}
\section{はさみうちの定理}

\begin{lem}[[C12]]
  $x, y, z, \alpha, \varepsilon \in \R$ について $x \leq z \to z \leq y \to \abs{x - \alpha} < \varepsilon \to \abs{y - \alpha} < \varepsilon \to \abs{z - \alpha} < \varepsilon$.
\end{lem}
\begin{proof}[通常の証明]
  [A11] より $-\varepsilon < x - \alpha$, $y - \alpha < \varepsilon$. $x - \alpha \leq z - \alpha \leq y - \alpha$ ([F19]) なので $-\varepsilon < z - \alpha < \varepsilon$ ([O3][O4]). [A12] より $\abs{z - \alpha} < \varepsilon$.
\end{proof}

\trees{ErrSqueeze}
\begin{thm}[はさみうちの定理. テキスト定理2.3]\label{thm:squeeze}
  $\ZFC \vdash \Seq(c) \to \forall n\,(n \in \N \to a_n \leq c_n \land c_n \leq b_n) \to \Conv(a, \alpha) \to \Conv(b, \alpha) \to \Conv(c, \alpha)$.
\end{thm}
\begin{proof}[通常の証明(テキストと同じ)]
  $\varepsilon > 0$ に対し補題\ref{lem:tail}で $N$ をとると, $n \geq N$ で $\abs{a_n - \alpha} < \varepsilon$, $\abs{b_n - \alpha} < \varepsilon$. $a_n \leq c_n \leq b_n$ と [C12] より $\abs{c_n - \alpha} < \varepsilon$.
\end{proof}

\trees{Squeeze}
\section{極限と大小関係}

\begin{thm}[テキスト定理2.5]\label{thm:lim-le}
  $\ZFC \vdash \Conv(a, \alpha) \to \Conv(b, \beta) \to \forall n\,(n \in \N \to a_n \leq b_n) \to \alpha \leq \beta$.
\end{thm}
\begin{proof}[通常の証明(テキストと同じ)]
  背理法で $\beta < \alpha$ と仮定し, $\varepsilon := (\alpha - \beta)/2 > 0$ とおく. 補題\ref{lem:tail}で $N$ をとると $\alpha - a_N < \varepsilon$, $b_N - \beta < \varepsilon$ ([A11], [F22], [F14]). また $a_N - b_N \leq 0$ なので $a_N - \beta = (a_N - b_N) + (b_N - \beta) \leq \varepsilon$. よって $\alpha - \beta = (\alpha - a_N) + (a_N - \beta) < \varepsilon + \varepsilon = \alpha - \beta$ となり矛盾. 従って $\neg\, \beta < \alpha$ であり, [O8] より $\alpha \leq \beta$.
\end{proof}

\trees{LimLe}

\section{数列の演算と有界性}

数列 $a, b$, 実数 $c$ に対して, 次の数列を定義 $(x \in A \mapsto t)$ (定義\ref{dfn:lambda})で定める. これらが数列であり値が右辺であることは基礎定理 [S2]--[S5], [S10]--[S12] である.
\begin{gather*}
  a \otimes b := (n \mapsto a_n b_n),\qquad a^{\ast} := (n \mapsto a_n^{-1})\ (\text{すべての } a_n \neq 0 \text{ のとき}),\\
  \overline{c} := (n \mapsto c),\qquad \abs{a} := (n \mapsto \abs{a_n}).
\end{gather*}
また, 数列が有界であることを $\mathrm{SeqBdd}(a) :\leftrightarrow \Seq(a) \land \exists M\,(M \in \R \land \forall n\,(n \in \N \to \abs{a_n} \leq M))$ で定める.

\begin{lem}\label{lem:seq-aux}
  \lemlist{ConvExt,FinBound,ErrProd}
\end{lem}
\begin{proof}[通常の証明]
  [C13] 値が一致する数列は収束先も一致する: 定義より明らか. [C14] $N$ に関する帰納法(定理\ref{thm:induction})で, $M_N := \abs{a_1} + \cdots + \abs{a_N}$ に相当する上界を作る. [C15] $a_n b_n - \alpha\beta = a_n(b_n - \beta) + \beta(a_n - \alpha)$ ([F18]) と三角不等式による.
\end{proof}

\begin{thm}[収束列の有界性. テキスト定理2.11]\label{thm:conv-bdd}
  収束する数列は有界である.
\end{thm}
\begin{proof}[通常の証明(テキストと同じ)]
  $\Conv(a, \alpha)$ とし, $\varepsilon = 1$ に対する $N$ をとる. $n \geq N$ ならば $\abs{a_n} \leq \abs{a_n - \alpha} + \abs{\alpha} < 1 + \abs{\alpha}$. $n \leq N$ の項は [C14] の $M$ で抑えられるので, $M + 1 + \abs{\alpha}$ が上界である.
\end{proof}
\fs{ConvBdd}

\begin{thm}[積の極限. テキスト定理2.7]\label{thm:lim-prod}
  $\Conv(a, \alpha)$, $\Conv(b, \beta)$ ならば $\Conv(a \otimes b, \alpha\beta)$.
\end{thm}
\begin{proof}[通常の証明(テキストと同じ)]
  定理\ref{thm:conv-bdd}より $\abs{a_n} \leq M$ ($M > 0$). $\varepsilon > 0$ に対して, 補題\ref{lem:tail}を十分小さい正数に適用すると, $n \geq N$ で $\abs{a_n b_n - \alpha\beta} \leq M\abs{b_n - \beta} + \abs{\beta}\abs{a_n - \alpha} < \varepsilon$ ([C15]).
\end{proof}
\fs{LimProd}

\begin{thm}[商の極限. テキスト定理2.9]\label{thm:lim-quot}
  $\Conv(b, \beta)$, $\beta \neq 0$, すべての $n$ で $b_n \neq 0$ ならば $\Conv(b^{\ast}, \beta^{-1})$. さらに $\Conv(a, \alpha)$ ならば $\Conv(a \otimes b^{\ast}, \alpha\beta^{-1})$.
\end{thm}
\begin{proof}[通常の証明(テキストと同じ)]
  $\varepsilon := \abs{\beta}/2$ に対する $N_1$ をとると, $n \geq N_1$ で $\abs{b_n} > \abs{\beta}/2$ ([E26]). よって $\abs{b_n^{-1} - \beta^{-1}} = \abs{\beta - b_n}\abs{b_n}^{-1}\abs{\beta}^{-1} < 2\abs{\beta}^{-2}\abs{b_n - \beta}$ ([E21]) であり, 右辺は $\varepsilon$ より小さくできる. 後半は定理\ref{thm:lim-prod}による.
\end{proof}
\fs{LimRecip}
\fs{LimQuot}

\section{部分列}

テキストの部分列(定義6.5)は $a_{n_1}, a_{n_2}, \ldots$ ($n_1 < n_2 < \cdots$)であり, 添字の列 $s = (n_k)_k$ を明示して次のように定める. $\mathrm{NSeq}(s)$ は $s : \N \to \N$ を表す.
\begin{dfn}[テキスト定義6.5]
  \begin{align*}
    \mathrm{SubIdx}(s) &:\leftrightarrow \mathrm{NSeq}(s) \land \forall k\,(k \in \N \to s_k < s_{S(k)}),\\
    \mathrm{SubSeq}(b, a) &:\leftrightarrow \Seq(b) \land \exists s\,\bigl(\mathrm{SubIdx}(s) \land \forall k\,(k \in \N \to b_k = a_{s_k})\bigr).
  \end{align*}
  $a \circ s$ は数列 $(k \mapsto a_{s_k})$ である([S12]).
\end{dfn}

\begin{lem}
  \lemlist{IdxLt,IdxGe,SubConvIdx}
\end{lem}
\begin{proof}[通常の証明]
  [C16] $m < n \to s_m < s_n$: $n$ に関する帰納法. [C17] $k \leq s_k$: $k$ に関する帰納法と [N10]. [C18] $\varepsilon$ に対する $N$ をとると, $k \geq N$ なら $s_k \geq k \geq N$ なので $\abs{a_{s_k} - \alpha} < \varepsilon$.
\end{proof}

\begin{thm}[収束する数列の部分列. テキスト定理6.7]\label{thm:subconv}
  $\Conv(a, \alpha)$ かつ $b$ が $a$ の部分列ならば $\Conv(b, \alpha)$.
\end{thm}
\begin{proof}[通常の証明]
  [C18] と [C13] による.
\end{proof}
\fs{SubConv}

偶数番目・奇数番目の添字列 $\mathsf{e} = (2k)_k$, $\mathsf{o} = (2k-1)_k$ について, これらが狭義単調で $\N$ を覆うことは基礎定理 [S15] である.

\begin{lem}\lemlist{CoverConv}\end{lem}
\begin{proof}[通常の証明]
  2つの添字列 $s, t$ が $\N$ を覆い, $a \circ s$, $a \circ t$ がともに $\alpha$ に収束するとする. $\varepsilon$ に対する $N_1, N_2$ をとり $N := \max\{s_{N_1}, t_{N_2}\}$ とおく. $n \geq N$ は $n = s_k$ または $n = t_k$ と書け, [C16] より $k \geq N_1$ (または $k \geq N_2$) なので $\abs{a_n - \alpha} < \varepsilon$.
\end{proof}

\begin{prop}[テキスト命題2.13]
  $\Conv(a, \alpha) \leftrightarrow \Seq(a) \land \Conv(a \circ \mathsf{o}, \alpha) \land \Conv(a \circ \mathsf{e}, \alpha)$.
\end{prop}
\begin{proof}[通常の証明]
  ($\Rightarrow$) 定理\ref{thm:subconv}. ($\Leftarrow$) [C19].
\end{proof}
\fs{OddEven}

\section{絶対値と発散}

\begin{lem}\lemlist{SubZero,AbsAbs,InvInv}\end{lem}

\begin{thm}[テキスト定理2.15]
  \begin{enumerate}[label=(\arabic*)]
    \item $\Conv(a, \alpha) \to \Conv(\abs{a}, \abs{\alpha})$.
    \item $\Seq(a) \to (\Conv(a, 0) \leftrightarrow \Conv(\abs{a}, 0))$.
    \item $\Seq(a) \to \alpha \in \R \to (\Conv(a, \alpha) \leftrightarrow \Conv(a \oplus \overline{-\alpha}, 0))$.
  \end{enumerate}
\end{thm}
\begin{proof}[通常の証明(テキストと同じ)]
  (1) $\bigl\lvert\abs{a_n} - \abs{\alpha}\bigr\rvert \leq \abs{a_n - \alpha}$ ([E27]). (2) $\bigl\lvert \abs{a_n} - 0 \bigr\rvert = \abs{a_n - 0}$ ([E28][E29]). (3) $(a_n + (-\alpha)) - 0 = a_n - \alpha$.
\end{proof}
\fs{LimAbs}
\fs{LimAbsZero}
\fs{LimSubConst}

テキストの $\lim a_n = \infty$ は $\mathrm{DvgP}(a)$ である(定義1.2(2)).

\begin{thm}[テキスト定理2.16, 2.17]
  \begin{enumerate}[label=(\arabic*)]
    \item $\mathrm{DvgP}(a)$ かつすべての $a_n \neq 0$ ならば $\Conv(a^{\ast}, 0)$.
    \item すべての $a_n > 0$ かつ $\Conv(a, 0)$ ならば $\mathrm{DvgP}(a^{\ast})$.
  \end{enumerate}
\end{thm}
\begin{proof}[通常の証明(テキストと同じ)]
  (1) $\varepsilon > 0$ に対し $K := \varepsilon^{-1}$ とすると, $n \geq N$ で $a_n > K > 0$ なので $\abs{a_n^{-1}} < \varepsilon$ ([E19]). (2) $K > 0$ に対し $\varepsilon := K^{-1}$ とすると, $n \geq N$ で $0 < a_n < K^{-1}$ なので $a_n^{-1} > K$ ([E19][E30]).
\end{proof}
\fs{DvgRecip}
\fs{RecipDvg}

\begin{thm}[テキスト定理2.18]
  $\mathrm{DvgP}(a)$ とする.
  \begin{enumerate}[label=(\arabic*)]
    \item すべての $n$ で $a_n \leq b_n$ ならば $\mathrm{DvgP}(b)$.
    \item すべての $n$ で $C \leq b_n$ ならば $\mathrm{DvgP}(a \oplus b)$.
    \item $C > 0$ かつすべての $n$ で $C \leq b_n$ ならば $\mathrm{DvgP}(a \otimes b)$.
  \end{enumerate}
\end{thm}
\begin{proof}[通常の証明(テキストと同じ)]
  $K > 0$ に対して, (1) は同じ $N$, (2) は $K - C$ (の正の部分)に対する $N$, (3) は $K C^{-1}$ に対する $N$ をとればよい.
\end{proof}
\fs{DvgLe}
\fs{DvgAdd}
\fs{DvgMul}

\section{等比数列の極限}

\begin{textref}{第3章 3.1節 命題3.1}
  等比数列 $r^n$ の極限.
\end{textref}

$r \in \R$ に対して, 累乗の列 $(r^n)_{n \in \N}$ を再帰定理(定理\ref{thm:recursion-omega})で $r^1 = r$, $r^{S(n)} = r^n \cdot r$ と定める. これが数列であることと再帰式は基礎定理 [S24] である. 本文中では $r^n$ と書き, 列全体を $r^{\bullet}$ と書く.

\begin{lem}\label{lem:pow}
  \lemlist{MulZeroCancel,PowAbs,PowNonneg,PowPos,PowNe0,PowOneOne,PowInv,PowDec,SucIdx}
\end{lem}
\begin{proof}[通常の証明]
  [E48] $x = x(yy^{-1}) = (xy)y^{-1} = 0$. [P1]--[P6] はいずれも $n$ に関する帰納法(第\ref{ch:naturals}章)であり, 帰納段階はそれぞれ [A17], [A16], [E11], [E15], (R8), [E16] による. [P7] $r^{S(n)} = r^n r \leq r^n \cdot 1$ ([E9]). [C28] $\mathsf{s} := (k \mapsto S(k))$ は添字列である([S14], [N17]).
\end{proof}

\begin{prop}[等比数列の極限. テキスト命題3.1]\label{prop:geo}
  $r \in \R$ とする.
  \begin{enumerate}[label=(\arabic*)]
    \item $r > 1$ ならば $\lim r^n = \infty$.
    \item $\lim 1^n = 1$.
    \item $\abs{r} < 1$ ならば $\lim r^n = 0$.
  \end{enumerate}
\end{prop}
\begin{proof}[通常の証明]
  (3) テキストは二項定理による評価を用いるが, ここでは単調収束定理(定理\ref{thm:monoconv})を用いる. $s := \abs{r}$ とおくと $0 \leq s < 1$ なので, $s^\bullet$ は単調減少で下に有界 ([P2][P7]) であり, ある $L$ に収束する. $s^\bullet \circ \mathsf{s} = (s^{n+1})_n$ は部分列なので $L$ に収束し([C18]), 一方 $s^{n+1} = s \cdot s^n$ なので $sL$ にも収束する(定理\ref{thm:lim-scalar}). 極限の一意性より $L = sL$, すなわち $L(1 - s) = 0$ であり, $1 - s \neq 0$ なので $L = 0$ ([E48]). $\abs{r^n} = s^n$ ([P1]) と定理2.15(2) より $\lim r^n = 0$.

  (2) [P5] と定数列の極限([G4]). (1) $q := r^{-1}$ とおくと $0 < q < 1$ ([E19], [E35]) なので (3) より $q^n \to 0$, $q^n > 0$ ([P3]). 定理2.17 より $(q^n)^{-1} \to \infty$ であり, $(q^n)^{-1} = r^n$ ([P6], [E30]) なので定理2.18(1) より $r^n \to \infty$.
\end{proof}
\fs{GeoLim}
\fs{GeoOne}
\fs{GeoDvg}

テキストの $r \leq -1$ の場合($\lim r^n$ は存在しない)は, 命題\ref{prop:geoser}(2) の証明と同様に $\abs{r^n} \geq 1$ から従うが, 形式化していない.

\section{級数}

\begin{textref}{第3章 4節 定義4.1, 命題4.3, 定理4.4, 定理4.7}
  級数の収束の定義, 無限等比級数, 収束する級数の和・定数倍, 級数が収束するための必要条件.
\end{textref}

数列 $a$ の部分和の列 $\Sigma a$ を再帰定理で $(\Sigma a)_1 = a_1$, $(\Sigma a)_{S(n)} = (\Sigma a)_n + a_{S(n)}$ と定める(基礎定理 [S25]). 本文中では $(\Sigma a)_n = \sum_{k=1}^n a_k$ と書く.

\begin{dfn}[テキスト定義4.1]
  級数 $\sum_{n=1}^\infty a_n$ が $s$ に収束するとは $\Conv(\Sigma a, s)$ のことである. このとき $\sum_{n=1}^\infty a_n = s$ と書く.
\end{dfn}

級数に関する定理では, 数列を1つずらすこと, および $\omega$ の 0 でない元が後者であること(定理\ref{thm:omega-order}(5). 基礎定理 [N18])を用いる.

\begin{lem}\label{lem:series}
  \lemlist{ShiftConv,AddSwap4,NegOneMul,PSumAdd,PSumScal,GeoPartial,PowGeOne}
\end{lem}
\begin{proof}[通常の証明]
  [C29] $(a_{n+1})_n$ が $\alpha$ に収束するとし, $\varepsilon$ に対する $N$ をとる. $n \geq S(N)$ ならば $n \neq 1$ なので $n = S(m)$ ([N18]) と書け, $S(N) \leq S(m)$ より $N \leq m$ ([N9], [O7]). よって $\abs{a_n - \alpha} = \abs{a_{S(m)} - \alpha} < \varepsilon$. [E49][E50] は体の公理から. [Q1][Q2][Q3][P8] は $n$ に関する帰納法である. 例えば [Q3] の帰納段階は, $p := r^{n+1}$, $q := \sum_{k=1}^n r^k$ として
  \[ (1 - r)(q + p) = (1-r)q + (1-r)p = (r - p) + (p - pr) = r - pr = r - r^{n+2}. \]
\end{proof}

\begin{thm}[収束する級数の和・定数倍. テキスト定理4.4]
  $\sum a_n = s$, $\sum b_n = t$, $\lambda \in \R$ ならば $\sum (a_n + b_n) = s + t$, $\sum \lambda a_n = \lambda s$.
\end{thm}
\begin{proof}[通常の証明]
  部分和について $\sum_{k=1}^n (a_k + b_k) = \sum_{k=1}^n a_k + \sum_{k=1}^n b_k$ ([Q1]), $\sum_{k=1}^n \lambda a_k = \lambda \sum_{k=1}^n a_k$ ([Q2]) なので, 定理\ref{thm:lim-sum}, \ref{thm:lim-scalar}と [C13] による. テキストの $\sum(\alpha a_n + \beta b_n) = \alpha s + \beta t$ はこの2つを組み合わせて得られる.
\end{proof}
\fs{SerAdd}
\fs{SerScal}

\begin{thm}[級数が収束する必要条件. テキスト定理4.7(1)]\label{thm:serterm}
  $\Seq(a)$ かつ $\sum a_n$ が収束するならば $\lim a_n = 0$.
\end{thm}
\begin{proof}[通常の証明(テキストと同じ)]
  $S_n := \sum_{k=1}^n a_k \to s$ とする. $(S_{n+1})_n$ も $s$ に収束する([C28][C18]) ので, $a_{n+1} = S_{n+1} - S_n \to s - s = 0$. [C29] より $a_n \to 0$.
\end{proof}
\fs{SerTerm}

テキストの定理4.7(2)「$a_n \not\to 0$ ならば $\sum a_n$ は発散する」は (1) の対偶である.

\begin{prop}[無限等比級数. テキスト命題4.3]\label{prop:geoser}
  $r \in \R$ とする.
  \begin{enumerate}[label=(\arabic*)]
    \item $\abs{r} < 1$ ならば $\sum_{n=1}^\infty r^n = r(1 - r)^{-1}$.
    \item $\abs{r} \geq 1$ ならば $\sum_{n=1}^\infty r^n$ は収束しない.
  \end{enumerate}
\end{prop}
\begin{proof}[通常の証明]
  (1) [Q3] より $\sum_{k=1}^n r^k = (1-r)^{-1}(r - r^{n+1})$. 命題\ref{prop:geo}(3) と [C18] より $r^{n+1} \to 0$ なので, 右辺は $(1-r)^{-1} r$ に収束する. (2) 収束するとすれば定理\ref{thm:serterm}より $r^n \to 0$, 従って $\abs{r}^n = \abs{r^n} \to 0$. しかし $\abs{r}^n \geq 1$ ([P8]) なので $\varepsilon = 1$ に対して矛盾する.
\end{proof}
\fs{GeoSer}
\fs{GeoSerDvg}

テキストの命題4.3 は初項 $a$ の級数 $\sum_{n=1}^\infty a r^{n-1} = a(1-r)^{-1}$ の形で述べられている. これは (1) と定理4.4 から $a + a r^{-1}\sum r^n$ ($r \neq 0$ のとき) として得られるが, その形での形式化はしていない.

\section{有界な単調数列}

\begin{textref}{第3章 5節 定義5.1, 定理5.3, 定理5.10}
  単調数列, 有界な単調数列の収束性, 区間縮小法.
\end{textref}

\begin{dfn}[テキスト定義5.1]
  $\mathrm{MonoInc}(a) :\leftrightarrow \Seq(a) \land \forall n\,(n \in \N \to a_n \leq a_{S(n)})$, \quad $\mathrm{MonoDec}(a) :\leftrightarrow \Seq(a) \land \forall n\,(n \in \N \to a_{S(n)} \leq a_n)$.
\end{dfn}

\begin{lem}\label{lem:mono}
  \lemlist{MonoIncLe,MonoDecLe,SubLtSwap,LtAddSwap,BandAbs,SubLtSelf,SupApprox,InfApprox}
\end{lem}
\begin{proof}[通常の証明]
  [C20][C21] $m \leq n \to a_m \leq a_n$ ($a_m \geq a_n$): $n$ に関する帰納法(第\ref{ch:naturals}章の構成手続き \texttt{monoBuild}). [E31]--[E34] は順序体の計算. [C22] $s = \sup E$, $e > 0$ ならば $s - e$ は上界でないので $s - e < x$ をみたす $x \in E$ がある. [C23] は下限について同様.
\end{proof}

\begin{thm}[有界な単調数列の収束性. テキスト定理5.3]\label{thm:monoconv}
  \begin{enumerate}[label=(\arabic*)]
    \item 単調増加で上に有界な数列は収束する.
    \item 単調減少で下に有界な数列は収束する.
  \end{enumerate}
\end{thm}
\begin{proof}[通常の証明(テキストと同じ)]
  (1) 値域 $E := \{a_n \mid n \in \N\}$ ([S16]) は空でなく上に有界なので, 上限 $s$ が存在する((R16)). $\varepsilon > 0$ に対して [C22] で $a_N > s - \varepsilon$ をとると, $n \geq N$ で $s - \varepsilon < a_N \leq a_n \leq s$ ([C20]) なので $\abs{a_n - s} < \varepsilon$ ([E33]). (2) 下限を用いて同様.
\end{proof}
\fs{MonoConv}
\fs{MonoConvDec}

\begin{thm}[区間縮小法. テキスト定理5.10]
  $a$ が単調増加, $b$ が単調減少で, すべての $n$ で $a_n \leq b_n$ ならば, すべての $n$ で $a_n \leq x \leq b_n$ をみたす $x \in \R$ が存在する.
\end{thm}
\begin{proof}[通常の証明(テキストと同じ)]
  $a_n \leq b_n \leq b_1$ なので $a$ は上に有界であり, 定理\ref{thm:monoconv}より $x := \lim a_n$ が存在する. $a_n \leq x$ は [C20] と定理\ref{thm:lim-le}から. また各 $m$ について $a_n \leq b_m$ ($n \geq m$ なら $a_n \leq b_n \leq b_m$, $n \leq m$ なら $a_n \leq a_m \leq b_m$) なので, 定理\ref{thm:lim-le}より $x \leq b_m$.
\end{proof}
\fs{NestedInt}

\section{Bolzano--Weierstrass の定理と Cauchy 列}

\begin{textref}{第3章 6節 定理6.8, 定義6.9, 補題6.10, 定理6.11}
  有界な数列は収束する部分列をもつ. 数列が収束することと Cauchy 列であることは同値である.
\end{textref}

部分列の抽出には次の基礎定理を用いる. これは「各 $k$ と各 $N$ について条件をみたす $n \geq N$ が存在する」ならば, 再帰定理で $s_{k+1}$ を「$s_k$ より大きく条件をみたす最小の $n$」と定めることで($\omega$ の整列性. 選択公理は不要)得られる.
\begin{itemize}
  \item[{[S28]}] \fsf{F:ExtractSub}
\end{itemize}

\begin{lem}\lemlist{ArchInv,IotaInvMono}\end{lem}
\begin{proof}[通常の証明]
  [C24] アルキメデスの公理(命題3.2)で $1 < \iota(K)\varepsilon$ となる $K$ をとる. [C25] [N16] と [E20].
\end{proof}

\begin{thm}[Bolzano--Weierstrass の定理. テキスト定理6.8]\label{thm:bw}
  $\mathrm{SeqBdd}(a)$ ならば, $a$ の部分列 $b$ で収束するものが存在する.
\end{thm}
\begin{proof}[通常の証明]
  テキストは区間縮小法による. ここでは上限を用いる. $\abs{a_n} \leq M$ とし, $E := \{x \in \R \mid x \leq M \land \forall N \in \N\,\exists n \geq N\,(x \leq a_n)\}$ とおく($x \leq a_n$ をみたす $n$ が無限にある $x$ の全体). $-M \in E$ であり $M$ は $E$ の上界なので $c := \sup E$ が存在する. 任意の $e > 0$ について, (i) $c + e \notin E$ なので $a_n \geq c + e$ となる $n$ は有限個, (ii) [C22] より $c - e < x$ となる $x \in E$ があるので $a_n > c - e$ となる $n$ は無限にある. よって任意の $k$, $N$ について $n \geq N$ かつ $\abs{a_n - c} < \iota(k)^{-1}$ をみたす $n$ がある. [S28] により $\abs{a_{s_k} - c} < \iota(k)^{-1}$ をみたす添字列 $s$ をとると, [C24][C25] より $a \circ s$ は $c$ に収束する.
\end{proof}
\fs{BW}

\begin{dfn}[Cauchy 列. テキスト定義6.9]
  \begin{align*}
    \mathrm{Cauchy}(a) :\leftrightarrow{}& \Seq(a) \land \forall\varepsilon\,\bigl(\varepsilon \in \R \to 0 < \varepsilon \to \exists N\,(N \in \N \\
    & \land \forall m\,\forall n\,(m, n \in \N \to N \leq m \to N \leq n \to \abs{a_m - a_n} < \varepsilon))\bigr).
  \end{align*}
\end{dfn}

\begin{lem}\lemlist{CauchyE,ErrChain}\end{lem}

\begin{lem}[Cauchy 列の有界性. テキスト補題6.10]
  Cauchy 列は有界である.
\end{lem}
\begin{proof}[通常の証明(テキストと同じ)]
  $\varepsilon = 1$ に対する $N$ をとると $n \geq N$ で $\abs{a_n} < \abs{a_N} + 1$. 残りの有限個は [C14] で抑える.
\end{proof}
\fs{CauchyBdd}

\begin{thm}[実数の完備性. テキスト定理6.11]
  $\mathrm{Cgt}(a) \leftrightarrow \mathrm{Cauchy}(a)$.
\end{thm}
\begin{proof}[通常の証明(テキストと同じ)]
  ($\Rightarrow$) $\abs{a_m - a_n} \leq \abs{a_m - \alpha} + \abs{\alpha - a_n}$. ($\Leftarrow$) 補題6.10 と定理\ref{thm:bw}より $a$ は $\alpha$ に収束する部分列 $a \circ s$ をもつ. $\varepsilon/2$ に対する Cauchy 列の $N_1$ と部分列の $N_2$ をとり, $k := \max\{N_1, N_2\}$ とすると, $n \geq N_1$ で $\abs{a_n - \alpha} \leq \abs{a_n - a_{s_k}} + \abs{a_{s_k} - \alpha} < \varepsilon$ ([C17], [C27]).
\end{proof}
\fs{CauchyComplete}

\section{本章で形式化していない内容}

テキスト第3章のうち, 次は形式化していない: 3.2--3.4節の計算例($\sqrt[n]{a} \to 1$, $n!$ と指数関数の比較など. 累乗根や階乗の定義が必要), 定理3.8, 5.3節の漸化式の例, 定義5.5(Napier 数 $e$)と関連する極限, 6.1節の公式(定理6.1, 命題6.2), 定理6.12(完備性とアルキメデスの公理から上限定理を導く. 本ノートでは上限定理を $\R$ の構成から直接証明しているので不要).
